\section*{第 \ref{ch:b2:meiosis-heredity} 章——减数分裂、遗传重组与遗传}

\begin{solution}{exo:b2:meiosis-heredity:1}
有丝分裂：一次分裂，同源染色体不配对，形成两个细胞，每个细胞
在遗传上都与亲代相同（$2n \to 2n$）。\omterm{def:b2:meiosis-heredity:meiosis}{减数分裂}：一次复制后
连续两次分裂，同源染色体配对并发生\omterm{prop:b2:meiosis-heredity:stages}{交叉互换}，形成四个细胞，
均为\omterm{def:b2:meiosis-heredity:meiosis}{单倍体}，且彼此遗传组成互不相同（$2n \to n$）。
\end{solution}

\begin{solution}{exo:b2:meiosis-heredity:2}
$2^{23} = 8.4\times 10^{6}$；$2^{4} = 16$；$2^{7} = 128$——在
\omterm{prop:b2:meiosis-heredity:stages}{交叉互换}之前，每一项都会无限制地成倍增加。
\end{solution}

\begin{solution}{exo:b2:meiosis-heredity:3}
$Tt\times Tt$：\omterm{def:b2:meiosis-heredity:vocabulary}{基因型}为 $1\,TT : 2\,Tt : 1\,tt$，\omterm{def:b2:meiosis-heredity:vocabulary}{表型}为 $3$ 高茎
：$1$ 矮茎。$Tt\times tt$：$1\,Tt : 1\,tt$，$1 : 1$。将高茎植株与
矮茎植株测交：若后代全为高茎，则为 $TT$；若一半为矮茎，则为 $Tt$
（或让其自交：$TT$ 稳定遗传，$Tt$ 按 $3 : 1$ 分离）。
\end{solution}

\begin{solution}{exo:b2:meiosis-heredity:4}
\omterm{def:b2:meiosis-heredity:linkage}{连锁}基因位于同一条染色体上，其亲本型
\omterm{def:b2:meiosis-heredity:vocabulary}{等位基因}组合在配子中占比过高；重组率
是重组配子所占的比例（即测交后代中的
重组型个体）；一\omterm{def:b2:meiosis-heredity:linkage}{厘摩}即重组率为百分之一。
位于长染色体上相距较远的基因，在几乎每次\omterm{def:b2:meiosis-heredity:meiosis}{减数分裂}中两者之间
都至少发生一次交换，多次交换使组合随机化，
因此 $r$ 达到 $\tfrac12$——
与独立分配无法区分。
\end{solution}

\begin{solution}{exo:b2:meiosis-heredity:5}
总数 556；期望值 $312.75 : 104.25 : 104.25 : 34.75$；$\chi^2 =
2.25^2/312.75 + 3.75^2/104.25 + 3.25^2/104.25 + 2.75^2/34.75 = 0.016
+ 0.135 + 0.101 + 0.218 = 0.47 < 7.81$：与 $9 : 3 : 3 :
1$ 相符。
\end{solution}

\begin{solution}{exo:b2:meiosis-heredity:6}
与 $1 : 1 : 1 : 1$（每类 250）相比：$\chi^2 = (162^2 + 138^2 + 147^2
+ 153^2)/250 = 361 \gg 7.81$：两基因\omterm{def:b2:meiosis-heredity:linkage}{连锁}。重组型 $103 + 97 = 200$
占 1000 的：$r = 0.20$，\qty{20}{cM}。在相斥组中，$Ab/aB$ 产生 400
个 $Ab$、400 个 $aB$、100 个 $AB$、100 个 $ab$：距离相同，亲本型
与重组型的类别互换。
\end{solution}

\begin{solution}{exo:b2:meiosis-heredity:7}
$d = -\tfrac12\ln(1 - 2r)$：$r = 0.20$ 时 $-\tfrac12\ln 0.6 =
0.255$，即 \qty{25.5}{cM}；$r = 0.45$ 时 $-\tfrac12\ln 0.1 = 1.15$，
即 \qty{115}{cM}。$r = \tfrac12(1 - e^{-2d})$：\qty{30}{cM} 时
$\tfrac12(1 - e^{-0.6}) = 0.226$；\qty{100}{cM} 时 $\tfrac12(1 -
e^{-2}) = 0.43$。在 \qty{10}{cM} 以下，校正可以忽略；
在 \qty{30}{cM} 以上，校正不可或缺；在 \qty{50}{cM} 以上，$r$
趋于饱和，单次两点测交已无法测出
距离——对远距离基因须通过近距离基因逐段连接来作图。
\end{solution}

\begin{solution}{exo:b2:meiosis-heredity:8}
(a) $X^wX^w \times X^+Y$：子代雄性全部白眼（$X^wY$），雌性全部
红眼（$X^+X^w$）。(b) $X^+X^w \times X^wY$：雄性一半红眼、一半白眼；
雌性一半红眼（$X^+X^w$），一半白眼（$X^wX^w$）。(c) 雌性
$X^+X^w$ $\times$ 兄弟 $X^wY$：与 (b) 相同——每种性别中
一半白眼、一半红眼。
\end{solution}

\begin{solution}{exo:b2:meiosis-heredity:9}
两种酶串联形成色素；紫色需要两个位点上都有
\omterm{def:b2:meiosis-heredity:vocabulary}{显性}\omterm{def:b2:meiosis-heredity:vocabulary}{等位基因}（$C\_P\_$）。两个纯系为 $CCpp$ 和 $ccPP$
（各自因不同原因呈白色）；\omterm{def:b2:meiosis-heredity:vocabulary}{杂合}体 $CcPp$ 为紫色；
自交后，紫色 $9\,C\_P\_$ 与白色 $7$ 之比（$3\,C\_pp + 3\,ccP\_ +
1\,ccpp$）。测交 $CcPp\times ccpp$：$1$ 紫色（$CcPp$）：$3$
白色。
\end{solution}

\begin{solution}{exo:b2:meiosis-heredity:10}
亲本型 $+++$、$abc$；双交换型 $+b+$、$a+c$；$b$ 是
位置发生互换的基因：顺序为 $a$--$b$--$c$。$r_{ab} = (11 + 9 + 2)/1202
= 0.018$（\qty{1.8}{cM}）；$r_{bc} = (118 + 112 + 2)/1202 = 0.193$
（\qty{19.3}{cM}）。双交换的期望数为 $0.018\times 0.193\times 1202 =
4.2$；观察值为 2：并发率 $0.47$，干涉 $0.53$。
\end{solution}

\begin{solution}{exo:b2:meiosis-heredity:11}
第二次分裂型子囊 $300/1000 = \qty{30}{\%}$；距离 $=
\tfrac12\times 30 = \qty{15}{cM}$。基因与
着丝粒之间的交换只涉及四条染色单体中的两条，因此
一个呈第二次分裂分离的子囊含两条重组型与两条
亲本型染色单体：重组率为此类子囊频率的
一半。对于 $2 : 4 : 2$：基因座与着丝粒之间的
交换会互换内侧的两条染色单体，因此减数第一次分裂时，
每个半纺锤体各带一条 $A$ 和一条 $a$
染色单体；减数第二次分裂随后按 $A\,a \mid
a\,A$（或相反顺序）分离它们，最后的有丝分裂使每个孢子数目加倍：
$2 : 4 : 2$。
\end{solution}

\begin{solution}{exo:b2:meiosis-heredity:12}
她的母亲是必然携带者（她从父亲那里接受了唯一的
X），因此该女子获得该\omterm{def:b2:meiosis-heredity:vocabulary}{等位基因}的概率为 $\tfrac12$。
两个健康儿子：携带者生出两个健康儿子的概率为 $(\tfrac12)^2 = \tfrac14$，
非携带者的概率为 1。由贝叶斯公式：
$P(\text{携带者}\mid\text{数据}) = (\tfrac12\times\tfrac14)/(\tfrac12
\times\tfrac14 + \tfrac12\times 1) = \tfrac{1/8}{5/8} = \tfrac15$.
\end{solution}

\begin{solution}{pb:b2:meiosis-heredity:1}
\textbf{1.} 亲本为 $vg/vg\;+/+$ 和 $+/+\;e/e$；$F_1$ 为 $+/vg\;+/e$；
配子 $+\,+$、$+\,e$、$vg\,+$、$vg\,e$，各占 $\tfrac14$。
\textbf{2.} $900 : 300 : 300 : 100$.
\textbf{3.} $\chi^2 = 8^2/900 + 9^2/300 + 4^2/300 + 3^2/100 = 0.07 +
0.27 + 0.05 + 0.09 = 0.48 < 7.81$：独立分配成立。
\textbf{4.} 两个位点均为\omterm{def:b2:meiosis-heredity:vocabulary}{纯合}野生型的个体占全部的 $\tfrac{1}{16}$，
野生型占 $\tfrac{9}{16}$：即 $\tfrac19$。
\textbf{5.} 野生型、残翅、黑檀体、残翅黑檀体，各占 $\tfrac14$。
\textbf{6.} \omterm{def:b2:meiosis-heredity:vocabulary}{表型}是\omterm{def:b2:meiosis-heredity:vocabulary}{基因型}与环境共同作用的产物
（温度影响色素合成酶）：\omterm{def:b2:meiosis-heredity:vocabulary}{基因型}规定的是
反应规范，而非固定的性状。
\textbf{7.} $b\,+/+\,vg$（相斥组）：$b$ 与 $+_{vg}$ 来自一个
亲本，$vg$ 与 $+_b$ 来自另一个亲本。
\textbf{8.} 亲本型：黑体（$b\,+$）965 和残翅（$+\,vg$）944；
重组型：黑体残翅（$b\,vg$）206 和野生型（$+\,+$）185。
野生型染色体 $+\,+$ 在 $F_1$ 中并不存在；它只能
通过\omterm{prop:b2:meiosis-heredity:stages}{交叉互换}产生。
\textbf{9.} $r = 391/2300 = 0.17$：\qty{17}{cM}。
\textbf{10.} $d = -\tfrac12\ln(1 - 0.34) = -\tfrac12\ln 0.66 = 0.21$:
\qty{21}{cM}.
\textbf{11.} 雄性\emph{Drosophila}中不发生\omterm{prop:b2:meiosis-heredity:stages}{交叉互换}：
雄配子只携带两条亲本染色体。
\textbf{12.} $F_1$ $b\,vg/+\,+$（偶联）：亲本型为野生型和
黑体残翅，各占 \qty{41.5}{\%}；重组型为黑体和
残翅，各占 \qty{8.5}{\%}。
\textbf{13.} $X^hX^+$：女儿从父亲那里接受唯一的 X，该 X
携带 $h$，并从未患病的母亲那里接受一条正常的 X。
\textbf{14.} $\tfrac12$（儿子）$\times$ $\tfrac12$（接受 $X^h$）$=
\tfrac14$.
\textbf{15.} $\tfrac12$：她从母亲的两条 X
染色体中任取一条；父亲提供正常的 X。
\textbf{16.} 每个孩子不患病的概率为 $\tfrac34$：
$(\tfrac34)^3 = \tfrac{27}{64} = 0.42$.
\textbf{17.} 儿子从父亲那里获得 Y，从不获得 X。女性
只有在 $X^hX^h$ 时才患病：需要患病的父亲与携带者（或
患病）母亲，这种情况很少见。
\textbf{18.} 先验概率 $\tfrac12$。两个健康儿子的概率
在她为携带者时为 $\tfrac14$，否则为 1：后验 $(\tfrac12\times
\tfrac14)/(\tfrac12\times\tfrac14 + \tfrac12) = \tfrac15$.
\textbf{19.} 儿子唯一的 X 来自母亲，而 Y 不携带
这些基因，因此每个儿子都直接显现出母方一个配子的
性状：本杂交本身即为测交。
\textbf{20.} 亲本型 $+++$（853）、$ywm$（833）；双交换型
$+w+$（2）、$y+m$（2）；$w$ 位置互换：顺序为 $y$--$w$--$m$。
\textbf{21.} $r_{yw} = (24 + 26 + 2 + 2)/2000 = 0.027$；$r_{wm} = (128
+ 132 + 2 + 2)/2000 = 0.132$。图谱：$y$——\qty{2.7}{cM}——$w$——
\qty{13.2}{cM}——$m$.
\textbf{22.} 期望值 $0.027\times 0.132\times 2000 = 7.1$；观察值
为 4：并发率 $0.56$，干涉 $0.44$。
\textbf{23.} $y$--$m$ 重组型：$y++$、$+wm$、$yw+$、$++m$：$(24 +
26 + 128 + 132)/2000 = 0.155$，对比 $0.027 + 0.132 = 0.159$：四类
双交换个体恢复了亲本的 $y$--$m$ 组合，因而
被漏计，$2\times 4/2000 = 0.004$。
\textbf{24.} 女儿从父亲那里接受带 $y\,w\,m$ 的 X，并
从母亲那里接受一条 X：所有女儿均为\omterm{def:b2:meiosis-heredity:vocabulary}{杂合}子，\omterm{def:b2:meiosis-heredity:vocabulary}{表型}为野生型，无论
其携带的是哪种重组型 X；重组因此不可见。
\textbf{25.} $y$——\qty{2.7}{cM}——$w$——\qty{13.2}{cM}——$m$;
并发系数 $0.56$，干涉 $0.44$。
\end{solution}
